% !TeX root = ../../Book3.tex
\section{调和共轭}
\label{b3:sec:1502}

在二维，寻找调和共轭等价于对一个闭微分形式积分。方程解决局部存在性，沿闭曲线的积分则决定局部对象能否拼成一个全局函数。

\subsection{Cauchy–Riemann 方程}
\label{b3:subsec:150201}

若希望 \(u+\mathrm{i}v\) 全纯，必须满足 \(v_x=-u_y\)、\(v_y=u_x\)。所需的一形式为 \(\omega=-u_y\,\mathrm{d}x+u_x\,\mathrm{d}y\)，其外微分为 \(\Delta u\,\mathrm{d}x\wedge\mathrm{d}y\)。所以调和方程恰好消去局部积分的路径依赖。这里的符号由 Cauchy--Riemann 方程确定，改成 \(u_y\,\mathrm{d}x-u_x\,\mathrm{d}y\) 会得到共轭的负号。

\subsection{局部调和共轭}
\label{b3:subsec:150202}

\begin{proposition}\label{b3:prop:15-local-conjugate}
实值调和函数在每点的某个圆盘邻域具有调和共轭，且共轭只差实常数。
\end{proposition}
\begin{proof}
以圆盘中心 \((a,b)\) 为起点，令
\[
v(x,y)=-\int_a^x u_y(t,b)\,\mathrm{d}t+\int_b^y u_x(x,t)\,\mathrm{d}t.
\]
得到 \(v_y=u_x\)，又由 \(u_{xx}=-u_{yy}\) 得 \(v_x=-u_y(x,b)-\int_b^yu_{yy}(x,t)\,\mathrm{d}t=-u_y(x,y)\)。因此 \(u+\mathrm{i}v\) 全纯，\(v\) 也调和。两个共轭的偏导相同，差在连通圆盘上为常数。
\end{proof}
例如 \(u=x^2-y^2\) 对应 \(v=2xy\)；\(u=\log|z|\) 的局部共轭为某个辐角分支。后一例在绕原点一周后增加 \(2\pi\)，已显示全局障碍。

\subsection{单连通区域}
\label{b3:subsec:150203}

\begin{theorem}\label{b3:thm:15-global-conjugate}
区域 \(D\) 上的实值调和函数 \(u\) 有单值调和共轭，当且仅当每条分段光滑闭曲线 \(\gamma\subset D\) 均满足 \(\int_\gamma\omega=0\)。特别地，单连通区域上一定存在共轭。
\end{theorem}
\begin{proof}
若 \(\omega=\mathrm{d}v\)，闭曲线积分为零。反之，从固定点到 \(z\) 积分 \(\omega\)；两条路径之差是闭曲线，故积分独立于路径，局部微分得到上一小节的偏导。单连通时，每条闭曲线可缩到一点；把同伦像用有限个局部原函数圆盘覆盖、细分参数方格，相邻边积分相消，故积分在同伦中不变。
\end{proof}
去心圆盘上的 \(\log|z|\) 给出 \(\omega=\mathrm{d}\arg z\)，单位圆上积分为 \(2\pi\)，所以单连通条件不能随意删去。

\subsection{全纯函数}
\label{b3:subsec:150204}

如果只需要全纯函数的模，可以先构造局部共轭。不同共轭的差为常数，故局部函数 \(e^{u+\mathrm{i}v}\) 在交集只差模一常数；它们不一定拼成单值函数。若所有周期落在 \(2\pi\mathbb Z\)，指数却能消除差异，得到全局无零点全纯函数。

\begin{proposition}\label{b3:prop:15-log-modulus}
区域上无零点全纯函数 \(f\) 的 \(\log|f|\) 调和。反之，调和函数 \(u\) 能写成 \(\log|f|\)，其中 \(f\) 无零点全纯，当且仅当 \(-u_y\,\mathrm{d}x+u_x\,\mathrm{d}y\) 的所有周期属于 \(2\pi\mathbb Z\)。
\end{proposition}
\begin{proof}
局部对数使第一结论成立，并使必要性化为辐角沿闭曲线的总变化。充分性由上面的路径积分定义 \(e^{u+\mathrm{i}\int\omega}\)，不同路径改变整数倍 \(2\pi\)，故指数相同；局部表示保证全纯。
\end{proof}


可直接从实部计算全纯函数的例子能帮助核对共轭符号。设
\(u=x^3-3xy^2+xy\)。先算 \(\Delta u=0\)，再由 \(v_y=u_x=3x^2-3y^2+y\) 积分得
\(v=3x^2y-y^3+y^2/2+h(x)\)。代入 \(v_x=-u_y=6xy-x\)，得到 \(h'(x)=-x\)，故
\[
v=3x^2y-y^3+\frac{y^2-x^2}{2}+C,
\quad f(z)=z^3-\frac{\mathrm{i}}2z^2+\mathrm{i}C.
\]
这个计算同时检查了调和条件、共轭的相容性和纯虚常数的唯一自由度。若一开始给出的 \(u\) 不调和，例如 \(x^2+y^2\)，两偏导方程会在混合导数处矛盾；这比形式地积分其中一条方程更能发现问题。

\ProseParagraphBreak

不同边界孔的周期可以各自贡献。在去掉 \(a_1,\ldots,a_k\) 的区域上，函数
\(u(z)=\sum c_j\log|z-a_j|\) 调和。沿只绕 \(a_j\) 一周的曲线，周期为 \(2\pi c_j\)。因此全局单值共轭存在当且仅当这些基本周期全部零；全局无零点全纯函数的对数模存在则要求 \(c_j\in\mathbb Z\)。后一情况由 \(f=\prod(z-a_j)^{c_j}\) 明确实现，在当前去点区域内负指数不会引入内部极点。两种存在问题的差异来自指数函数的周期，不是来自不同的调和方程。

\ProseParagraphBreak

单连通性是保证所有调和函数都有共轭的条件，却不是某个特定函数有共轭的必要条件。环域上 \(u=\operatorname{Re}z\) 有共轭 \(\operatorname{Im}z\)，而 \(\log|z|\) 没有。对一个给定 \(u\)，应检查它的闭形式周期；不能只看到域有孔就拒绝任何共轭构造。


\ProseParagraphBreak

局部共轭还使调和函数具有实解析性。在小圆盘写 \(u=\operatorname{Re}f\)，全纯函数的 Taylor 级数局部一致收敛，因此
\[
u(a+re^{\mathrm{i}\theta})=c_0+\sum_{n\ge1}r^n(A_n\cos n\theta+B_n\sin n\theta).
\]
这里 \(c_0\) 为实数，系数由 \(f\) 的复 Taylor 系数拆成实部、虚部得到。每个非恒定项本身调和，且圆周平均零。若 \(u\) 在某个非空开集恒零，则其局部全纯函数实部恒零，Cauchy--Riemann 方程使它为纯虚常数；沿重叠圆盘传播后，\(u\) 在连通区域恒零。这个唯一延拓结论比最大值原理要求的单侧极值条件不同：前者使用开集上的全部零值，后者使用一点极值。

\ProseParagraphBreak

若两个调和函数在一条曲线上相等，唯一延拓却未必适用。\(u=x\) 与零函数在虚轴相等但在两侧不同。若额外给出该直线上的法向导数相等，在具有解析延拓的局部情形，Taylor 系数才能被同时限制。只由曲线上的函数值不能恢复二维的全部局部数据；全纯函数沿有聚点的零集恒等定理不能直接挪给实值调和函数。

\ProseParagraphBreak

Harnack 不等式还把局部正性传遍连通区域。固定紧集 \(K\Subset D\) 和一点 \(z_0\in D\)，用有限条圆盘链把 \(z_0\) 连到覆盖 \(K\) 的小圆盘；每一步比较的常数只依赖圆盘的半径比例。相乘后得到 \(\sup_Ku\le C_{K,z_0}u(z_0)\)，对所有正调和函数使用同一个常数。连通性在这里转化为有限次局部比较，而非直接给一个全域常数。若区域无界或紧集逐渐碰到边界，常数一般随之增大；圆盘 Poisson 核靠近其边界质量所在点时的增长就是一个具体检验。
